\subsection{ERD hệ thống OLTP và luồng dữ liệu}
Sơ đồ quan hệ thực thể của hệ thống bán hàng có \code{orders} làm trung tâm: mỗi khách (\code{customers}) đặt nhiều đơn; mỗi đơn có nhiều dòng hàng (\code{order\_\allowbreak{}items}), nhiều lần thanh toán (\code{order\_\allowbreak{}payments}) và nhiều đánh giá (\code{order\_\allowbreak{}reviews}). Dòng hàng nối tới sản phẩm (\code{products}) và người bán (\code{sellers}). Toàn bộ nguồn có 11 khoá chính và 7 khoá ngoại. Khoá chính của \code{order\_\allowbreak{}items} là cặp (\code{order\_\allowbreak{}id}, \code{order\_\allowbreak{}item\_\allowbreak{}id}), của \code{order\_\allowbreak{}payments} là (\code{order\_\allowbreak{}id}, \code{payment\_\allowbreak{}sequential}). Với \code{order\_\allowbreak{}reviews}, mã \code{review\_\allowbreak{}id} xuất hiện lặp ở nhiều đơn nên không dùng làm khoá chính được; khoá là cặp (\code{review\_\allowbreak{}id}, \code{order\_\allowbreak{}id}). Bảng \code{geolocation} không có khoá tự nhiên, vì một mã bưu chính có nhiều toạ độ, nên được thêm cột kỹ thuật \code{id}; nó chỉ liên hệ logic với khách và người bán qua mã bưu chính. Cũng chỉ liên hệ logic là \code{products} và bảng dịch tên danh mục, vì có danh mục không có bản dịch. Hệ thống CRM gồm \code{mql} và \code{closed\_\allowbreak{}deals} nối theo quan hệ 1–0..1 (một khách tiềm năng có thể được chốt thành một thương vụ). Cột \code{closed\_\allowbreak{}deals.seller\_\allowbreak{}id} nối logic sang \code{oltp.sellers} bằng đường nét đứt trong sơ đồ và không có khoá ngoại. Mã bưu chính trong bản dữ liệu này đã đầy đủ 5 chữ số, kể cả số 0 đầu, nên được lưu dạng văn bản.

\hinh{3-2-4_erd_oltp.png}{ERD hai hệ thống nguồn: bán hàng (oltp) và marketing (crm)}{fig:erd-oltp}{1.0}
\todo{Yêu cầu của môn: trên ERD OLTP \textbf{gạch chân các trường dữ liệu sẽ dùng sang OLAP} và ghi chú \textbf{mapping} tương ứng. Vẽ lại ERD bằng draw.io từ hình nháp, gạch chân các cột dùng và ghi chú (bảng mapping đầy đủ ở Phụ lục C, nguồn docs/mapping\_oltp\_olap.md). Thêm sơ đồ luồng dữ liệu (data flow) từ nguồn đến kho.}
